\section{把芯片写成寄存器与标志位结构}

端口号只是地址，不能代表完整设备。要理解一颗芯片，至少要同时记录寄存器宽度、
读写方向、字段位、复用条件和状态所有权。项目把当前启动链使用的 CPU、16550A
UART 和 ATA/IDE 主通道整理为机器可读的 \filepath{docs/hardware/register_map.yaml}，
本节提炼最常用的走读视图。

\subsection{CPU 标志与模式控制}

\begin{longtable}{L{0.12\textwidth}L{0.16\textwidth}L{0.56\textwidth}}
\toprule
\textbf{位} & \textbf{名称} & \textbf{当前含义}\\
\midrule
\endfirsthead
\toprule
\textbf{位} & \textbf{名称} & \textbf{当前含义}\\
\midrule
\endhead
0 & CF & 无符号加法进位；地址和 LBA 加法后用于溢出判断。\\
2 & PF & 低字节偶校验；当前不作为协议条件。\\
6 & ZF & 结果为零；由 \code{cmp}、\code{test} 和循环条件消费。\\
8 & TF & 单步陷阱；正常启动保持关闭。\\
9 & IF & 可屏蔽中断开关；固件入口用 \code{CLI} 关闭。\\
10 & DF & 字符串方向；入口用 \code{CLD} 清零，保证 \code{lodsb/lodsw} 向高地址走。\\
11 & OF & 有符号溢出；当前不作为硬件协议条件。\\
\bottomrule
\end{longtable}

模式转换还会写 CR0.PE、CR0.PG、CR3、CR4.PAE 与 IA32\_EFER.LME。CR0.PG
不能孤立打开：页表根、PAE、EFER 和远跳转顺序共同决定处理器最终状态。
EFER.LMA 是处理器激活后的只读结果，软件不能把“请求写入”误认为“已经进入长模式”。

\subsection{描述符寄存器与异常位字段}

v0.5 把 GDTR、IDTR、TR 和 CR2 加入结构化硬件图。GDTR/IDTR 都由 64 位
base 与 16 位 inclusive limit 组成；TR 保存 TSS 选择子和处理器内部缓存。
v0.8 当前 GDT limit=55，IDT limit=4095，TR=\code{0x28}。

\begin{longtable}{L{0.10\textwidth}L{0.18\textwidth}L{0.22\textwidth}L{0.34\textwidth}}
\toprule
\textbf{对象} & \textbf{宽度} & \textbf{访问指令} & \textbf{当前职责}\\
\midrule
\endfirsthead
\toprule
\textbf{对象} & \textbf{宽度} & \textbf{访问指令} & \textbf{当前职责}\\
\midrule
\endhead
GDTR & 80 位 & LGDT/SGDT & 七槽 GDT 的 base/limit\\
IDTR & 80 位 & LIDT/SIDT & 256 槽 IDT 的 base/limit\\
TR & 16 位可见 & LTR/STR & 选择 104 字节 TSS\\
CR2 & 64 位 & MOV & 最近页故障线性地址\\
\bottomrule
\end{longtable}

页故障错误码 bit 0 是 present/permission，bit 1 是 read/write，bit 2 是
supervisor/user，bit 3 是 reserved-bit violation，bit 4 是
data/instruction fetch。Ring 0 not-present 注入结果为全零，CR2 为
\code{0x04000000}；Ring 3 读取未映射页得到 \code{0x4}，其 U/S 位证明
访问来自用户态。

机器可读 YAML 同时列出 32 个 present 异常门、16 个硬件 IRQ 门、
\code{0x80} 系统调用门、十个硬件错误码向量，以及
NMI→IST2、双重故障→IST1、机器检查→IST3 的结构关系。这样芯片文档、C++
位字段测试和 NASM 异常桩共享同一组可审查事实。

\subsection{16550A UART 的复用寄存器}

COM1 基址为 \code{0x3F8}，寄存器宽度为 8 位。LCR 的 DLAB 位使偏移 0 和 1
在发送/接收寄存器与除数锁存器之间复用：

\begin{longtable}{L{0.12\textwidth}L{0.30\textwidth}L{0.38\textwidth}}
\toprule
\textbf{偏移} & \textbf{名称} & \textbf{关键字段}\\
\midrule
\endfirsthead
\toprule
\textbf{偏移} & \textbf{名称} & \textbf{关键字段}\\
\midrule
\endhead
0 & RBR/THR/DLL & DLAB=0 时接收或发送；DLAB=1 时除数低字节。\\
1 & IER/DLM & DLAB=0 时中断使能；DLAB=1 时除数高字节。\\
2 & IIR/FCR & 读中断识别，写 FIFO 控制。\\
3 & LCR & WLS、停止位、奇偶校验和 DLAB。\\
4 & MCR & DTR、RTS、OUT2 等 modem 控制。\\
5 & LSR & DR、错误位、THRE、TEMT 和 FIFO 错误。\\
\bottomrule
\end{longtable}

LSR bit 5 的 THRE 只表示发送保持寄存器可接收新字节；bit 6 的 TEMT 表示保持
寄存器和移位寄存器都空，语义更强但不是本阶段写入的必要条件。固件因此只等待
THRE，并给每次等待设置 \code{0xFFFF} 次预算。

\subsection{ATA/IDE 状态位决定数据所有权}

主通道命令块端口为 \code{0x1F0..0x1F7}，控制端口为 \code{0x3F6}。DATA
是 16 位寄存器，每个 512 字节扇区必须读取 256 个字；LBA 地址分散在三个
8 位端口和 DRIVE/HEAD 的四个高位中。

\begin{longtable}{L{0.10\textwidth}L{0.18\textwidth}L{0.56\textwidth}}
\toprule
\textbf{位} & \textbf{名称} & \textbf{决策}\\
\midrule
\endfirsthead
\toprule
\textbf{位} & \textbf{名称} & \textbf{决策}\\
\midrule
\endhead
0 & ERR & 设备报告命令错误，立即进入 \code{IDE\_ERROR}。\\
3 & DRQ & 数据可读；必须同时满足 BSY=0 且 ERR/DF=0。\\
5 & DF & 设备故障，与 ERR 归入同一设备错误边界。\\
6 & DRDY & 设备就绪，当前不单独推进状态。\\
7 & BSY & 设备忙，继续有界等待；预算耗尽进入 \code{IDE\_TIMEOUT}。\\
\bottomrule
\end{longtable}

因此一条读取事务的状态机是：

\begin{center}
\code{BSY=1}\ $\rightarrow$ 等待；\quad
\code{ERR/DF=1}\ $\rightarrow$ 设备错误；\quad
\code{BSY=0,DRQ=1}\ $\rightarrow$ 读取 DATA。
\end{center}

把这些字段写成结构后，驱动策略就不再散落魔法数字。将来加入 PIC、PIT、PS/2、
PCI 或 LAPIC 时，先扩展芯片寄存器描述，再实现访问层和状态机，最后让策略层只
依赖命名后的状态。
